\documentclass[10pt,a4paper]{amsart}
\usepackage[margin=19mm]{geometry}
\usepackage{amsmath,amssymb,mathtools}
\usepackage[scr=rsfs,cal=euler]{mathalfa}
\usepackage{xfrac}
\usepackage[unicode,hidelinks,bookmarks=false]{hyperref}
\newcommand{\ie}{i.e.\ }
\newcommand{\cf}{cf.\ }
\newcommand{\resp}{respectively\ }
\newcommand{\Subsection}{\subsection}
\providecommand{\nobreakdash}{\nobreak\hbox{-}}
\setlength{\emergencystretch}{2em}
\multlinegap0pt
\usepackage{amssymb}
\usepackage[shortlabels]{enumitem}
\newtheorem{proposition}{Proposition}
\newtheorem{theorem}{Theorem}
\newcommand{\JT}{JT}
\newcommand{\norm}[1]{\left\lVert#1\right\rVert}
\newcommand{\ran}{\operatorname{ran}}
\newcommand{\supp}{\operatorname{supp}}
\title{A Study of the Extreme Points in the Unit Ball of $\JT$}
\author{Spiros A. Argyros}
\date{Source: arXiv:2602.23207v2 (CC BY 4.0)}
\begin{document}
\maketitle

\noindent\emph{Source and licence.} A Study of the Extreme Points in the Unit Ball of $\JT$. Version arXiv:2602.23207v2 (CC BY 4.0), distributed by arXiv under the Creative Commons Attribution 4.0 International licence, http://creativecommons.org/licenses/by/4.0/, as recorded in the arXiv licence metadata for this exact version and shown on the abstract page https://arxiv.org/abs/2602.23207v2. Full licence notice: \texttt{licenses/arxiv-2602.23207-CC-BY-4.0.txt}.\par\smallskip

\noindent\textit{Editorial scope.} Counted unit: the article's theorem theorem (source0.tex lines 944--946). The article's complete original proof of it is delivered with it (source0.tex lines 948--1002, one \texttt{proof} environment of 55 lines). No mathematics is written, repaired or re-derived: the statement and the proof are the article's own lines, in the article's own notation, and no retained line is substituted or re-flowed.  Retained uncounted as the source-local context the proof needs: lines 232--239; lines 295--297; lines 840--842.  Nothing was omitted from any retained span, so the package carries no omission record.  No retained line contains a citation, so the wrapper carries no bibliography and no bibliography entry.  No retained line contains a figure environment, an image-inclusion command or drawing code, so no image is delivered and the package declares no asset dependency.  Editorial formatting only, in addition: an AMS \texttt{amsart} preamble that restates the article's own declarations and macros used by the retained lines, namely the environments \texttt{proposition}, \texttt{theorem} on separate counters, together with the standard packages those lines need.  The retained environments print in this document as Proposition 1, Theorem 1 in order of appearance, whereas the article prints the same items with the numbering of its own full document.  The complete native source of the article as distributed in the arXiv e-print for this exact version is bundled unchanged as \texttt{sources/195385/source0.tex}.
\begin{proposition}[Characterization of Extreme Points]\label{prop.1}
    A vector $x \in \JT$ is \textbf{not} an extreme point if and only if there exists a non-zero vector $y \in \JT$ such that:
    \begin{enumerate}
        \item $\norm{x}_{\JT} = \norm{x+y}_{\JT} = \norm{x-y}_{\JT}$, and
        \item For every $x$-norming partition $\mathcal{P} = \{S_i\}$, we have $\norm{y}_{\mathcal{P}} = 0$.
    \end{enumerate}
    Note that $\norm{y}_{\mathcal{P}} = 0$ implies that for every segment $S \in \mathcal{P}$, $\sum_{\alpha \in S} y(\alpha) = 0$.
\end{proposition}

\begin{proposition}\label{prop:not_sep_not_ext}
    Let $x \in \JT$. If $x$ is not separated, then $x$ is not an extreme point.
\end{proposition}

\begin{theorem}\label{T5.9}
    Let $x \in \JT$ be a positive vector. The partition $\mathcal{S}$ generated by the Greedy Algorithm is a canonical $x$-norming partition.
\end{theorem}

\begin{theorem}
    Let $x \in \JT$ be a positive vector ($x \ge 0$). Then $x$ is an extreme point if and only if $x$ is separated.
\end{theorem}

\begin{proof}
    \textbf{($\Rightarrow$) Direction:}
    We have already established in Proposition \ref{prop:not_sep_not_ext} that if a vector is not separated, it cannot be an extreme point. This holds for all vectors, including positive ones.
    
    \textbf{($\Leftarrow$) Direction:}
    Assume that $x$ is separated but, for the sake of contradiction, that $x$ is \textbf{not} an extreme point.
    
    By the characterization of non-extreme points (Proposition \ref{prop.1}), there exists a non-zero vector $y \in \JT$
    with $\supp(y) \subset \ran(x)$ such that:
    \begin{enumerate}
        \item $\norm{x}_{\JT} = \norm{x+y}_{\JT} = \norm{x-y}_{\JT}$.
        \item For every $x$-norming partition $\mathcal{P}$, and every segment $S \in \mathcal{P}$, the sum of the perturbation vanishes: $\sum_{\gamma \in S} y(\gamma) = 0$.
    \end{enumerate}
    
    Since $y \neq 0$, we can choose a node $a \in \supp(y)$. Observe that $a \in \ran(x)$.
    
    \textbf{Step 1: Isolate $a$ from above.}
    Since $x$ is separated, there exists an $x$-norming partition $\mathcal{P}_1$ that separates $a$ from its parent (if $a \neq \emptyset$). Let $s_a \in \mathcal{P}_1$ be the segment containing $a$. Because $a$ is separated from its parent, $a$ must be the minimal element (the starting node) of $s_a$.
    By the property of $y$, we must have:
    \begin{equation}\label{eq:sum_sa}
        \sum_{\gamma \in s_a} y(\gamma) = 0.
    \end{equation}
    
    If $s_a = \{a\}$, then $y(a) = 0$, which contradicts $a \in \supp(y)$. Thus, $s_a$ must contain at least one descendant of $a$ with respect to $\supp(x)$. Let $\beta$ be the child of $a$ contained in $s_a$. We can write $s_a = s_{a, \beta} \cup t_\beta$, where $t_\beta$ is a segment starting at $\beta$.
    Equation (\ref{eq:sum_sa}) becomes:
    \[ y(a) + \sum_{\gamma \in t_\beta} y(\gamma) = 0 \implies \sum_{\gamma \in t_\beta} y(\gamma) = -y(a) \neq 0. \]
    
        
    \textbf{Step 2: Isolate $a$ from below.}
    Since $x$ is separated, there exists another $x$-norming partition $\mathcal{P}_2$ that separates $a$ from $\beta$.
    In this partition, $a$ belongs to some segment $s' \in \mathcal{P}_2$ which does \textbf{not} contain $\beta$.
    This implies that the restriction of $\mathcal{P}_2$ to the subtree $W_\beta$ (rooted at $\beta$) is an $x_\beta$-norming partition. Let us denote the square value of the optimal partition on $W_\beta$ as $\norm{x_\beta}_{\JT}^2$.
    
    \textbf{Step 3: The Contradiction.}
    We return to the segment $s_a = s_{a, \beta} \cup t_\beta$ from Step 1.
    Since $\mathcal{P}_1$ is an $x$-norming partition, it must be consistent with the Greedy Algorithm. This implies that the tail $t_\beta$ is an optimal path starting at $\beta$. Specifically, $t_\beta$ must be part of some $x_\beta$-norming partition for the subtree $W_\beta$ (Theorem \ref{T5.9}).
    
    Let $\mathcal{P}_\beta$ be an $x_\beta$-norming partition that contains the segment $t_\beta$. We can construct a new partition $\mathcal{P}^*$ for the whole vector $x$ by combining:
    \begin{itemize}
        \item The partition $\mathcal{P}_\beta$ for the subtree $W_\beta$ (containing $t_\beta$).
        \item All $s \in \mathcal{P}_2$ such that $s \cap W_\beta = \emptyset$.
    \end{itemize}
    This partition $\mathcal{P}^*$ is $x$-norming.
    
    However, strictly speaking, we rely on the property of non-extreme points: \textbf{Any} $x$-norming partition must have $y$-sum zero on its segments.
    
    The segment $t_\beta$ is a segment in the valid $x$-norming partition $\mathcal{P}^*$.
    Thus, we must have:
    \[ \sum_{\gamma \in t_\beta} y(\gamma) = 0. \]
    
    But from Step 1, we established that:
    \[ \sum_{\gamma \in t_\beta} y(\gamma) = -y(a) \neq 0. \]
    
    This is a contradiction. Therefore, the assumption that $x$ is not an extreme point must be false.
\end{proof}

\end{document}
